% Livro sintético de demonstração do scan (D48): um livro-texto inteiro, com sumário, páginas
% iniciais em romano (o número impresso não é o índice do PDF), cabeçalho corrido, partes,
% capítulos, seções, exemplos, blocos de exercícios, itens, um exercício que vira a página,
% fórmulas e uma figura. Conteúdo original, livre para versionar.
\documentclass[11pt,oneside]{book}
% Preâmbulo comum das fixtures do scan (D48). Fontes Latin Modern com T1, para que a camada de
% texto traga os acentos como caracteres, e babel em português, para que \chapter imprima
% "Capítulo" como um livro brasileiro imprime.
\usepackage[T1]{fontenc}
\usepackage[utf8]{inputenc}
\usepackage[brazil]{babel}
\usepackage{lmodern}
\usepackage{amsmath,amssymb}
\usepackage{tikz}
\usepackage[a4paper,margin=2.5cm]{geometry}
\usepackage{enumitem}
\setlist[enumerate,2]{label=\alph*)}
\makeatletter\@ifundefined{c@chapter}{\newcounter{exemplo}}{\newcounter{exemplo}[chapter]}\makeatother
\newenvironment{exemplo}{\par\medskip\refstepcounter{exemplo}\noindent\textbf{Exemplo \theexemplo.}\ }{\par\medskip}
\newcommand{\blocoexercicios}[1][Exercícios]{\section*{#1}}

\usepackage{tikz}
\newcommand{\texto}{Este parágrafo desenvolve a teoria da seção com calma. Ele apresenta as
definições, comenta as consequências imediatas e prepara o leitor para os exemplos resolvidos
que vêm a seguir, sempre com a notação usual dos livros de matemática do ensino superior.\par}
\begin{document}
\frontmatter
\begin{titlepage}\centering\vspace*{5cm}
{\Huge\bfseries Matemática Sintética}\\[1cm]{\Large Volume único}\\[4cm]{\large Editora Fixture}
\end{titlepage}
\tableofcontents
\mainmatter
\part{Fundamentos}
\chapter{Números reais}
\section{Ordem}
\texto\texto
\begin{exemplo}
Como $2 < 3$ e $3 < 5$, a transitividade garante que $2 < 5$.
\end{exemplo}
\texto
\subsection{Valor absoluto}
O valor absoluto de $x$ é $|x| = \max\{x, -x\}$, e vale a desigualdade triangular
\[ |x + y| \le |x| + |y|. \]
\begin{exemplo}
$|-7| = 7$ e $|3 - 10| = 7$.
\end{exemplo}
\blocoexercicios
\begin{enumerate}[label=\textbf{\arabic*.}]
\item Prove que $|xy| = |x|\,|y|$.
\item Resolva as inequações:
  \begin{enumerate}
  \item $|x - 1| < 2$;
  \item $|2x + 3| \ge 5$.
  \end{enumerate}
\item Mostre que $\sqrt{2}$ é irracional.
\end{enumerate}
\section{Supremo}
\texto\texto\texto
\begin{exemplo}
O supremo de $\{1 - \frac{1}{n} : n \ge 1\}$ é $1$.
\end{exemplo}
\texto\texto\texto\texto\texto
\blocoexercicios[Exercícios propostos]
\begin{enumerate}[label=\textbf{\arabic*.}]
\item Determine o supremo e o ínfimo de $\{\frac{n}{n+1} : n \ge 1\}$.
\item \texto\texto\texto\texto\texto\texto Este exercício é longo de propósito, para virar a página: prove que todo
conjunto não vazio e limitado superiormente de números reais tem supremo, e responda:
  \begin{enumerate}
  \item o supremo pertence sempre ao conjunto?
  \item o que muda se o conjunto for finito?
  \item dê um exemplo em que o supremo não é o máximo.
  \end{enumerate}
\item Calcule $\sup\,\{x \in \mathbb{Q} : x^2 < 2\}$.
\end{enumerate}
\part{Cálculo}
\chapter{Derivadas}
\section{Definição}
\texto
A derivada de $f$ no ponto $a$ é
\[ f'(a) = \lim_{h\to 0} \frac{f(a+h) - f(a)}{h}. \]
\begin{center}\begin{tikzpicture}[scale=0.9]
\draw[->] (-0.2,0) -- (4,0); \draw[->] (0,-0.2) -- (0,3);
\draw[thick,domain=0:3.2] plot (\x,{0.25*\x*\x});
\draw (1,0.25) -- (3,2.25);
\end{tikzpicture}\end{center}
\texto
\begin{exemplo}
Se $f(x) = x^2$, então $f'(x) = 2x$.
\end{exemplo}
\section{Regras de derivação}
\texto\texto
\begin{exemplo}
Pela regra do produto, $(x \sin x)' = \sin x + x\cos x$.
\end{exemplo}
\blocoexercicios
\begin{enumerate}[label=\textbf{\arabic*.}]
\item Derive $f(x) = \dfrac{x^2 + 1}{x - 1}$.
\item Derive $g(x) = \sqrt{1 + x^2}$.
\item Calcule a derivada de $\det\begin{pmatrix} x & 1 \\ 2 & x \end{pmatrix}$.
\end{enumerate}
\end{document}
